% Компилировать pdfLaTeX
\documentclass[12pt,a4paper]{article}

\usepackage{cmap}
\usepackage[T2A]{fontenc}
\usepackage[utf8]{inputenc}
\usepackage[english,russian]{babel}
\usepackage{tempora}

\usepackage{amsmath,amssymb,mathtools}

\usepackage{geometry}
\geometry{left=25mm,right=15mm,top=20mm,bottom=20mm}

\usepackage{setspace}
\onehalfspacing
\usepackage{indentfirst}
\setlength{\parindent}{1.25cm}
\setlength{\parskip}{3pt}
\allowdisplaybreaks

\begin{document}

\begin{center}
    \textbf{\large Ч.\,2, №10}
\end{center}

\noindent
\textbf{Постановка:}
\[
J(x) = \int_0^\pi (\dot{x}^2 - x^2) dt \to \text{extr}
\]
\[
x(0) = 0, \qquad \int_0^\pi x \cos t \, dt = \frac{\pi}{2}, \qquad \int_0^\pi x \sin t \, dt = -2
\]

\noindent
\textbf{Замена \(\dot{x} = u\), введение \(y_1, y_2\):}
\[
y_1(t) = \int_0^t x \cos \tau \, d\tau, \qquad y_2(t) = \int_0^t x \sin \tau \, d\tau
\]
\[
\mathcal{L} = \int_0^\pi \bigl[ \lambda_0(u^2 - x^2) + p(\dot{x} - u) + q_1(\dot{y}_1 - x \cos t) + q_2(\dot{y}_2 - x \sin t) \bigr] dt + l
\]
\[
l = \lambda_1 x(0) + \lambda_2 y_1(0) + \lambda_3 y_2(0) + \lambda_4\left(y_1(\pi) - \frac{\pi}{2}\right) + \lambda_5(y_2(\pi) + 2)
\]

\noindent
\textbf{Функция Гамильтона:}
\[
H = pu + q_1 x \cos t + q_2 x \sin t - \lambda_0(u^2 - x^2)
\]

\vspace{0.3cm}
\noindent
\textbf{1. Стационарность по \(u\) (условие максимума ПМП):}
\[
\frac{\partial H}{\partial u} = p - 2\lambda_0 u = 0 \quad \Rightarrow \quad p = 2\lambda_0 u
\]
\[
\frac{\partial^2 H}{\partial u^2} = -2\lambda_0 \le 0 \quad \text{(для минимума)}
\]

\noindent
\textbf{2. Стационарность по переменным (\(x, y_1, y_2\)):}
\begin{itemize}
    \item \textbf{По \(x\):}
    \[
    \frac{\partial H}{\partial x} = q_1 \cos t + q_2 \sin t + 2\lambda_0 x, \quad \dot{p} = -\frac{\partial H}{\partial x} \quad \Rightarrow \quad \dot{p} = -2\lambda_0 x - q_1 \cos t - q_2 \sin t
    \]
    \item \textbf{По \(y_1\):}
    \[
    \frac{\partial H}{\partial y_1} = 0, \quad \dot{q}_1 = -\frac{\partial H}{\partial y_1} \quad \Rightarrow \quad \dot{q}_1 = 0 \quad \Rightarrow \quad q_1 = \text{const}
    \]
    \item \textbf{По \(y_2\):}
    \[
    \frac{\partial H}{\partial y_2} = 0, \quad \dot{q}_2 = -\frac{\partial H}{\partial y_2} \quad \Rightarrow \quad \dot{q}_2 = 0 \quad \Rightarrow \quad q_2 = \text{const}
    \]
\end{itemize}

\noindent
\textbf{3. Условия трансверсальности:}
\[
p(0) = \frac{\partial l}{\partial x(0)} = \lambda_1, \qquad p(\pi) = -\frac{\partial l}{\partial x(\pi)} = 0
\]
\[
q_1(0) = \frac{\partial l}{\partial y_1(0)} = \lambda_2, \qquad q_1(\pi) = -\frac{\partial l}{\partial y_1(\pi)} = -\lambda_4
\]
\[
q_2(0) = \frac{\partial l}{\partial y_2(0)} = \lambda_3, \qquad q_2(\pi) = -\frac{\partial l}{\partial y_2(\pi)} = -\lambda_5
\]

\noindent
\textbf{4. Нормальность и знаки:}
\[
\lambda_0 \ge 0 \ \text{(loc min)}
\]
При \(\lambda_0=0 \Rightarrow p=0 \Rightarrow q_1=q_2=0 \Rightarrow \lambda_1=\dots=\lambda_5=0\) (тривиально). Следовательно, \(\lambda_0 = 1\).
\[
p = 2u = 2\dot{x}
\]

\vspace{0.3cm}
\noindent
\textbf{5. Краевая задача:}
\[
\begin{cases}
\dot{x} = u = \frac{p}{2}, \\[1mm]
\dot{p} = -2x - q_1 \cos t - q_2 \sin t, \\[1mm]
\dot{q}_1 = 0, \\[1mm]
\dot{q}_2 = 0,
\end{cases}
\qquad
\begin{cases}
x(0) = 0, \\[1mm]
p(\pi) = 0, \\[1mm]
q_1(\pi) = -\lambda_4, \\[1mm]
q_2(\pi) = -\lambda_5.
\end{cases}
\]

Исключая \(p\) (\(\dot{p} = 2\ddot{x}\)):
\[
2\ddot{x} = -2x - q_1 \cos t - q_2 \sin t \quad \Rightarrow \quad \ddot{x} + x = -\frac{q_1}{2} \cos t - \frac{q_2}{2} \sin t
\]

\noindent
\textbf{6. Решение краевой задачи:}
\[
x(t) = C_1 \cos t + C_2 \sin t + a t \sin t + b t \cos t
\]
\[
\ddot{x} + x = 2a \cos t - 2b \sin t \quad \Rightarrow \quad a = -\frac{q_1}{4}, \quad b = \frac{q_2}{4}
\]
\[
x(0) = 0 \Rightarrow C_1 = 0
\]
\[
\dot{x}(t) = C_2 \cos t + a(\sin t + t \cos t) + b(\cos t - t \sin t)
\]
\[
p(\pi) = 2\dot{x}(\pi) = 0 \Rightarrow -C_2 - \pi a - b = 0 \Rightarrow C_2 = -\pi a - b
\]

\noindent
\textbf{7. Интегральные условия:}
\[
\int_0^\pi t \sin t \cos t \, dt = -\frac{\pi}{4}, \quad \int_0^\pi t \cos^2 t \, dt = \frac{\pi^2}{4}, \quad \int_0^\pi \sin t \cos t \, dt = 0
\]
\[
\int_0^\pi \sin^2 t \, dt = \frac{\pi}{2}, \quad \int_0^\pi t \sin^2 t \, dt = \frac{\pi^2}{4}
\]
\[
\int_0^\pi x \cos t \, dt = a\left(-\frac{\pi}{4}\right) + b\left(\frac{\pi^2}{4}\right) = \frac{\pi}{2}
\]
\[
\int_0^\pi x \sin t \, dt = C_2\left(\frac{\pi}{2}\right) + a\left(\frac{\pi^2}{4}\right) + b\left(-\frac{\pi}{4}\right) = -2
\]

\noindent
\textbf{8. Решение системы:}
\[
\begin{cases} 
-\dfrac{\pi}{4}a + \dfrac{\pi^2}{4}b = \dfrac{\pi}{2} \\[2mm]
C_2\dfrac{\pi}{2} + \dfrac{\pi^2}{4}a - \dfrac{\pi}{4}b = -2 \\[2mm]
C_2 + \pi a + b = 0 
\end{cases}
\Rightarrow
\begin{cases} 
-a + \pi b = 2 \\[1mm]
2\pi C_2 + \pi^2 a - \pi b = -8 \\[1mm]
C_2 = -\pi a - b 
\end{cases}
\]
\[
-\pi^2 a - 3\pi b = -8 \quad \Rightarrow \quad -\pi^3 b + 2\pi^2 - 3\pi b = -8
\]
\[
\pi(\pi^2 + 3)b = 2(\pi^2 + 4) \Rightarrow b = \frac{2(\pi^2 + 4)}{\pi(\pi^2 + 3)}
\]
\[
a = \pi b - 2 = \frac{2}{\pi^2 + 3}, \qquad C_2 = -\pi a - b = -\frac{4(\pi^2 + 2)}{\pi(\pi^2 + 3)}
\]

\noindent
\textbf{9. Экстремаль:}
\[
\hat{x}(t) = \frac{2}{\pi(3+\pi^2)} \left[ -2(2+\pi^2)\sin t + t\left((4+\pi^2)\cos t + \pi\sin t\right) \right]
\]

\vspace{0.3cm}
\noindent
\textbf{10. Глобальный минимум:}
\[
x = \hat{x} + h, \quad h(0) = 0, \quad \int_0^\pi h \cos t \, dt = \int_0^\pi h \sin t \, dt = 0
\]
\[
J(\hat{x}+h) - J(\hat{x}) = \int_0^\pi (2\dot{\hat{x}}\dot{h} - 2\hat{x}h) dt + \int_0^\pi (\dot{h}^2 - h^2) dt
\]
\[
2\int_0^\pi \dot{\hat{x}}\dot{h} \, dt = -2\int_0^\pi \ddot{\hat{x}}h \, dt \Rightarrow -2\int_0^\pi (\ddot{\hat{x}} + \hat{x})h \, dt = \int_0^\pi (q_1 \cos t + q_2 \sin t)h \, dt = 0
\]
\[
\Delta J = \int_0^\pi (\dot{h}^2 - h^2) dt, \quad v = \dot{h} \perp \{\sin t, 1+\cos t\}
\]
\[
\int_0^\pi h^2 dt \le C \int_0^\pi v^2 dt, \quad C = \frac{5\pi^4-84\pi^2+384}{6(3\pi^2-16)} \approx 0.5144 < 1
\]
\[
\Delta J \ge (1-C)\int_0^\pi v^2 dt > 0 \quad (v \not\equiv 0) \Rightarrow \hat{x} \text{ — строгий глобальный минимум}
\]

\vspace{0.3cm}
\noindent
\textbf{11. Отсутствие максимума:}
\[
H^{\max} = u^2 - x^2 + pu + \dots \ \Rightarrow \ H^{\max}_{uu} = 2 > 0 \ \Rightarrow \ \sup_{u} H^{\max} = +\infty
\]
\[
h = \alpha \sin 3t \ \Rightarrow \ h(0) = 0, \quad \int_0^\pi h \cos t \, dt = 0, \quad \int_0^\pi h \sin t \, dt = 0
\]
\[
\Delta J = \int_0^\pi (9\alpha^2\cos^2 3t - \alpha^2\sin^2 3t) dt = \alpha^2 \left( 9 \cdot \frac{\pi}{2} - \frac{\pi}{2} \right) = 4\pi\alpha^2
\]
\[
\lim_{|\alpha|\to\infty} J(\hat{x}+h) = J_{\min} + 4\pi\alpha^2 = +\infty \ \Rightarrow \ \text{глобального максимума нет.}
\]

\vspace{0.3cm}
\noindent
\textbf{12. Значение минимума:}
\[
J(\hat{x}) = \int_0^\pi (\dot{\hat{x}}^2 - \hat{x}^2) dt = [\hat{x}\dot{\hat{x}}]_0^\pi - \int_0^\pi \hat{x}(\ddot{\hat{x}} + \hat{x}) dt = 0 - \int_0^\pi \hat{x}\left(-\frac{q_1}{2}\cos t - \frac{q_2}{2}\sin t\right) dt
\]
\[
J(\hat{x}) = \frac{q_1}{2} \int_0^\pi \hat{x}\cos t \, dt + \frac{q_2}{2} \int_0^\pi \hat{x}\sin t \, dt = \frac{q_1}{2} \cdot \frac{\pi}{2} + \frac{q_2}{2} \cdot (-2) = \frac{\pi q_1}{4} - q_2
\]
Подставляя \(q_1 = -4a\) и \(q_2 = 4b\):
\[
J_{\min} = -\pi a - 4b = -\pi\left(\frac{2}{\pi^2+3}\right) - 4\left(\frac{2(\pi^2+4)}{\pi(\pi^2+3)}\right) = -\frac{2(5\pi^2+16)}{\pi(\pi^2+3)}
\]

\vspace{0.5cm}
%------------------------------------------------
\section*{Ответ}
%------------------------------------------------

\noindent
\[
\boxed{
\hat{x}(t) = \frac{2}{\pi(3+\pi^2)} \left[ -2(2+\pi^2)\sin t + t\left((4+\pi^2)\cos t + \pi\sin t\right) \right]
}
\]

\noindent
\[
\boxed{
\hat{u}(t) = \frac{2}{\pi(3+\pi^2)} \left[ -\pi^2\cos t + \pi\sin t + t\left(\pi\cos t - (4+\pi^2)\sin t\right) \right]
}
\]

\noindent
\[
\boxed{
J_{\min} = -\frac{2(5\pi^2+16)}{\pi(3+\pi^2)} \approx -3.2326
}
\]

\noindent
\textbf{Характер экстремума:}
\begin{itemize}
    \item \(\hat{x}(t)\) доставляет \textbf{строгий глобальный минимум} (доказано через положительность второй вариации: \(\Delta J \ge (1-C)\int \dot{h}^2 dt > 0\)).
    \item \textbf{Глобального максимума не существует}, так как для задачи на максимум функция Гамильтона выпукла по \(u\) (\(H^{\max}_{uu} = 2 > 0\)), и при вариации \(h(t) = \alpha \sin 3t\) имеем \(J \to +\infty\) при \(|\alpha| \to \infty\).
\end{itemize}

\end{document}